% Part III-B, sections that do not depend on the results (included by main_en.tex).

\section{Setting: a verifier that labels only some completions}\label{sec:setting}

Part II gives the verifier's verdict to every completion. Here only a fraction $f=0.25$ of each group's completions is
labelled, two of eight, chosen at random. The labelled share models a verifier that is costly to run, such as human
grading or slow tests. The question is how to use the unlabelled 75\%:
\begin{itemize}
\item \texttt{rlvr\_sub}: leave them out. The advantages come from the two labelled completions of each group, and the
  others get no gradient.
\item \texttt{sub\_kev}: reward them with \kev{}'s $\pc$ (``Is the final answer of the proposed solution correct?''), as
  released and frozen.
\item \texttt{sub\_kev\_recal}: reward them with $\pc$ from a copy of \kev{}'s pointer head that is refitted online on
  all labelled completions seen so far.
\item \texttt{sub\_probe}: reward them with a head of the same shape, on the same features, trained from scratch on the
  same labels. This separates what \kev{}'s released head contributes from what any small head fitted to these labels
  would.
\end{itemize}
Labelled completions always get the verifier's reward. The upper reference is Part II's \texttt{rlvr}, where every
completion is labelled. Everything else is Part II's setup: the same model, data, learning rate, batch, number of steps
and evaluation (\S2 of Part II). For logging only, the verifier also checks every unlabelled completion, which measures
how good each fill-in reward was on the completions it was used for.

\section{Online recalibration}\label{sec:recal}

\kev{}'s pointer head turns the backbone's hidden states into option logits $z$. Its $p(\text{yes})$ is
$\mathrm{softmax}(z/T)_{\text{yes}}$ with the released temperature $T$. \texttt{sub\_kev\_recal} keeps the backbone and
\kev{}'s adapter frozen and refits a float32 copy $\theta$ of the head after every step. The loss is the binary
cross-entropy of $p(\text{yes})$ against the verifier's labels, plus $\alpha\lVert\theta-\theta_0\rVert^2$, an anchor to
the released head $\theta_0$ ($\alpha=10^{-3}$). The fit uses 50 Adam steps (learning rate $10^{-4}$, batches of 256)
over a buffer of the last 8192 labelled completions. The features (the head's inputs) are computed once, when \kev{}
scores a completion, so the refit costs no backbone pass. The released head on the backbone is never changed.
\texttt{sub\_probe} uses the same procedure from a fresh initialisation of the head's projections, with $T=1$ and no
anchor.

The reason to expect recalibration to matter is that the policy's completions drift during training. They get longer
or shorter and change style, so a judge calibrated on its own training distribution can lose calibration on them. The
labelled completions measure that drift directly, and online refitting is the cheapest way to follow it.
